\textit{Physique I, année 2023 -- filière PC}

On cherche à déterminer à partir de quel moment on atteint une densité moyenne $n = 10^{11}\text{ m}^{-3}$ d'électrons caractéristique des sylphes. On note $\mathscr{N}$ le nombre d'étapes produisant des électrons lors des chocs. Pour simplifier, on considère que tous les électrons produits jusqu'à l'étape $\mathscr{N}$ sont contenus dans un volume dont la taille caractéristique est $\mathscr{N} \ell_{pm}$.

\begin{itemize}
    \item[\textbf{Q -- 24.}] Déterminer la valeur numérique de $\mathscr{N}$. Pour cette question, on prendra un libre parcours moyen $\ell_{pm}$ de l'ordre du centimètre. On pourra utiliser les courbes de la figure 7 en expliquant de façon détaillée la démarche suivie.
\end{itemize}

\TLEFigure{7}{Représentation graphique des fonctions $\mathscr{N} \mapsto 10^{-4} \times \mathscr{N}^3$ et $\mathscr{N} \mapsto 10^{-9} \times 2^{\mathscr{N}-3}$}
